% TOP_PROBLEM: {"id": 3700024, "title": "Goldberg's constant", "category_id": 8, "category": {"id": 8, "name": "analysis", "display_name": "Analysis", "description": "Limits, continuity, calculus, and function theory.", "slug": "analysis", "order_index": 8, "created_at": "2026-07-31T15:26:25.671Z"}, "status": "partially_solved", "problem_number": "AMR-036-0024", "set_id": 13, "statusQualification": "Fixes curvature -1 for the hyperbolic-radius convention; other metric conventions rescale this radius constant."}
% FIRST_POSTED: 2026-10-01
% ENTRY_KIND: statement_only
% UnsolvedMath ID 3700024; source catalogue code AMR-036-0024
% !TeX program = lualatex
% Definition and sources only; this file is not an LLM solution attempt.
\documentclass[11pt,a4paper]{article}
\usepackage[margin=25mm]{geometry}
\usepackage{fontspec}
\setmainfont{lmroman10-regular.otf}[BoldFont=lmroman10-bold.otf,
 ItalicFont=lmroman10-italic.otf,BoldItalicFont=lmroman10-bolditalic.otf]
\usepackage{amsmath,amssymb,mathtools}
\usepackage{unicode-math}
\setmathfont{latinmodern-math.otf}
\AtBeginDocument{\let\setminus\smallsetminus}
\usepackage{xurl,hyperref}
\hypersetup{colorlinks=true,urlcolor=blue}
\setcounter{secnumdepth}{0}
\setlength{\parindent}{0pt}
\setlength{\parskip}{5pt}
\setlength{\emergencystretch}{3em}
\begin{document}
\section*{Goldberg's constant}\label{3700024}
{\small UnsolvedMath ID 3700024 · AMR-036-0024 · Analysis · AMR Open Problem Lists}
\subsection{Definitions and mathematical statement}
Let $\mathbb D=\{z:|z|<1\}$ carry the hyperbolic metric
of curvature $-1$, with line element
$2|dz|/(1-|z|^2)$. Its distance is
\[
 d_{\mathbb D}(z,w)=2\operatorname{artanh}
                 \left|\frac{z-w}{1-\overline w z}\right|.
\]
Consider all holomorphic functions $f$ on $\mathbb D$
having exactly one zero $z_0$, of multiplicity one, and
exactly two $1$-points $z_1,z_2$, counted with multiplicity.
The two $1$-points may coincide, in which case that
single point is a zero of multiplicity two of $f-1$.

Define the hyperbolic enclosing radius of these marked
points by
\[
 R(f)=\inf_{c\in\mathbb D}
            \max_{j=0,1,2}d_{\mathbb D}(c,z_j),
 \qquad A_2=\inf_f R(f).
\]
Determine $A_2$, including sharpness and the functions
attaining the infimum if it is attained. Equivalently,
find the largest universal lower bound for the radius
of a hyperbolic disk containing all the zero and
$1$-point locations.

The associated structural question is whether every
extremal function must have a double $1$-point.
Disk automorphisms may move the locations without changing
their hyperbolic enclosing radius. No boundedness or
univalence of $f$ is assumed, but no additional zeros
or $1$-points anywhere in $\mathbb D$ are allowed.
The metric normalization is fixed here so that numerical
values of a radius constant are unambiguous.
\subsection{Short English statement}
Find the smallest hyperbolic disk that can contain the unique simple zero and two one-points of a holomorphic disk function.
\subsection{Sources}
\begin{itemize}
\item[S1] ULAM.AI and the UnsolvedMath Contributors, supplied problems.json, record 3700024 (AMR-036-0024), AMR Open Problem Lists. Catalogue identity and source transcription; CC BY 4.0.
\url{https://huggingface.co/datasets/ulamai/UnsolvedMath}.
\item[S2] Eremenko - My favorite unsolved problems, Goldberg's constant. Original source indexed by AMR-036-0024.
\url{https://www.math.purdue.edu/~eremenko/uns1.html}.
\item[S3] A. Eremenko, Goldberg's constant and its relatives, introductory problem. Author-maintained primary problem note.
\url{https://www.math.purdue.edu/~eremenko/dvi/goldconst.pdf}.
\end{itemize}
\end{document}
